\begin{afterword}
\summaryhead

The first post of the author's memoir sequence tells how the young Yudkowsky came to value intelligence above everything. The parents answered the child's arguments, including arguments against Judaism, by appeal to experience. The child concluded that people who downplay intelligence do not understand it. The older author corrects this: an eleven-year-old could not have out-thought adult parents in a fair contest, and the parents lost because they used their intelligence against itself.

The young author wanted everyone made smarter, claimed no moral superiority, and avoided the obvious traps of science-fiction villains. Then, in 1996, came the idea of the Singularity, and a ``happy death spiral'' around intelligence, shown by false beliefs such as that the speed of light would not limit a superintelligence. The real wrong turn was answering the objection that intelligence has nothing to do with goodness: the young author concluded that superintelligence would necessarily imply supermorality. A postscript says what has changed since.

\responsehead

The post is memoir, and it is judged as memoir: its reports of the young author's feelings and reasons are taken as given. What can be checked is its account of the young author's beliefs against what the young author wrote at the time.

The central report holds up. The archived text of ``Staring into the Singularity,'' dated 1996 and 1999, says: ``I am reasonably (95\%) certain that the Powers will be ethical,'' and ``If wiping out your creators is morally wrong, it won't happen.'' ``The Singularitarian Principles'' (2000, revised 2001) looks back on ``a point where I was sure that superintelligent meant super-ethical (probably true).'' The two texts differ in strength: the essay gives 95 per cent, and the later look back says ``sure.'' Its parenthesis, ``(probably true),'' is the verdict on the same claim when that text was written, so the belief was still held as probable in 2000 or 2001. The post's example of an earlier false belief, that the speed of light would be no barrier, I could not find stated, though ``Staring into the Singularity'' shows the attitude behind it.

The same essay bears on the young author's modesty. It does call its author dumb, as the post says the young author was careful to do. It also calls its author ``an actual enhanced human such as myself,'' which sits beside the post's ``No line drawn between himself and others.'' The essay is later than the 1995 the post describes, and neither passage claims moral superiority.

The label ``happy death spiral'' comes from Book II, and the notes on ``Affective Death Spirals'' discuss the evidence for its feedback step. Here the post offers a real case. ``Resist the Happy Death Spiral'' dropped its one realistic example of a false nice claim about a Great Idea; this post gives one that a careful admirer actually made, and it has the form of the halo effect: one good trait, intelligence, raising the perception of another, goodness.

One paragraph is ambiguous. It calls intelligence ``the unfairest'' card in the present tense, without saying whose view that is, and the postscript later puts death and old age first.

In short: a memoir whose central report, that the young author expected superintelligence to be moral, is borne out by the archived writings; its example of the spiral's earlier false beliefs, about the speed of light, I could not trace.
\end{afterword}
